% =========================================================================
% Capítulo 1: Cinemática
% Tema 01: Coordenadas cilíndricas
% Problema 14: Movimiento en Espiral Logarítmica
% =========================================================================
\subsubsection{Movimiento en Espiral Logarítmica}

Un punto $P$ sigue una trayectoria en espiral logarítmica descrita en coordenadas cilíndricas por:
$$r(\theta) = r_0 e^{-k\theta}$$

donde $r_0$ es el radio inicial, $k$ es una constante positiva de amortiguamiento, y $\theta$ es el ángulo azimutal. El punto se desplaza de modo que su velocidad angular es constante: $\dot{\theta} = \omega_0$.

\begin{enumerate}
    \item Expresar la velocidad del punto en coordenadas cilíndricas.
    \item Calcular la aceleración del punto.
    \item Demostrar que la rapidez (magnitud de velocidad) decrece exponencialmente con el tiempo.
\end{enumerate}

\begin{figure}[!ht]
    \centering
    \begin{tikzpicture}[>=Latex, scale=1.6]

        % --- CONFIGURACIÓN DE PARÁMETROS ---
        \def\ro{2.5}      % Radio inicial r0 [cm]
        \def\k{0.15}      % Constante de amortiguamiento k
        \def\w{1}         % Escala de la espiral

        % --- SOPORTES Y ELEMENTOS FIJOS ---
        % Origen
        \coordinate (O) at (0,0);
        \filldraw (O) circle (1.5pt) node[above left] {$O$};

        % --- GEOMETRÍA PRINCIPAL ---
        % Espiral logarítmica (dibujada como varios puntos conectados)
        \draw[thick, azul]
            plot[parametric, samples=100, domain=0:3*pi]
            ({0.8*\ro*exp(-\k*\x/pi)*cos(\x r)}, {0.8*\ro*exp(-\k*\x/pi)*sin(\x r)});

        % --- SISTEMAS DE COORDENADAS ---
        % Ejes de referencia
        \draw[->, thick] (O) -- (3, 0) node[right, font=\small] {$x$};
        \draw[->, thick] (O) -- (0, 3) node[above, font=\small] {$y$};

        % --- PUNTOS Y COORDENADAS ---
        % Punto inicial en la espiral (t=0)
        \coordinate (P0) at (0.8*\ro, 0);
        \filldraw[fill=white, thick] (P0) circle (2.5pt);
        \filldraw (P0) circle (1pt) node[right, font=\small] {$P(t=0)$};

        % Punto en tiempo intermedio
        \coordinate (Pm) at ({0.8*\ro*exp(-\k*pi/2)*cos(pi/2 r)}, {0.8*\ro*exp(-\k*pi/2)*sin(pi/2 r)});
        \filldraw[fill=white, thick] (Pm) circle (2pt);
        \filldraw (Pm) circle (1pt);

        % Punto cerca del origen (t grande)
        \coordinate (Pf) at ({0.8*\ro*exp(-\k*3*pi)*cos(3*pi r)}, {0.8*\ro*exp(-\k*3*pi)*sin(3*pi r)});
        \filldraw[fill=white, thick] (Pf) circle (2pt);
        \filldraw (Pf) circle (1pt) node[left, font=\small] {$P(t \to \infty)$};

        % --- VECTORES Y FUERZAS ---
        % Vector de velocidad en el punto inicial
        % Velocidad radial (hacia adentro) + velocidad tangencial
        \draw[->, thick, myred] (P0) -- ($(P0) + (-0.4, 0)$) node[below, font=\small, color=myred] {$\vec{v}_r$};
        \draw[->, thick, orange] (P0) -- ($(P0) + (0, 0.5)$) node[above, font=\small, color=orange] {$\vec{v}_\theta$};

        % Vector velocidad total (resultante)
        \draw[->, thick, gris1] (P0) -- ($(P0) + (-0.4, 0.5)$) node[above right, font=\small, color=gris1] {$\vec{v}$};

        % --- ANOTACIONES Y ETIQUETAS ---
        % Ángulo theta en un punto
        \draw[dashed, gray] (0, 0) -- (Pm);
        \draw[->, dashed, gray] (0.5, 0) arc (0:90:0.5);
        \node at (0.6, 0.4) [font=\small] {$\theta$};

        % Radio en el punto inicial
        \draw[<->, thin, gray] (0, -0.35) -- (P0) node[midway, below, font=\small] {$r_0$};

        % Etiqueta
        \node[font=\footnotesize] at (1.5, -2) {Espiral logarítmica: $r(\theta) = r_0 e^{-k\theta}$};

    \end{tikzpicture}
    \caption{Trayectoria en espiral logarítmica con velocidad descompuesta en componentes radial y transversal}
    \label{fig:c1_e01_p05_espiral}
\end{figure}

% --- SOLUCIÓN ---------------------------------------------------------------
\begin{solucion}{Espiral logarítmica}
\subsubsection*{Pregunta 1: Velocidad en Coordenadas Cilíndricas}

Dado que $\theta(t) = \omega_0 t$ (velocidad angular constante), tenemos $\dot{\theta} = \omega_0$ y $\ddot{\theta} = 0$.

De la ecuación de la espiral $r(\theta) = r_0 e^{-k\theta}$, obtenemos $r(t) = r_0 e^{-k\omega_0 t}$.

Derivando respecto al tiempo:
$$\dot{r} = r_0 \cdot (-k\omega_0) e^{-k\omega_0 t} = -k\omega_0 \cdot r_0 e^{-k\omega_0 t} = -k\omega_0 \cdot r(t)$$

La velocidad en coordenadas cilíndricas es:
$$\vec{v} = \dot{r} \, \hat{e}_r + r\dot{\theta} \, \hat{e}_\theta$$

Sustituyendo:
\begin{equation}
\boxed{\vec{v} = -k\omega_0 r(t) \, \hat{e}_r + r(t) \omega_0 \, \hat{e}_\theta = \omega_0 r(t) \left[ -k \, \hat{e}_r + \hat{e}_\theta \right]}
\end{equation}

\subsubsection*{Pregunta 2: Aceleración del Punto}

La aceleración en coordenadas cilíndricas es:
$$\vec{a} = (\ddot{r} - r\dot{\theta}^2) \hat{e}_r + (r\ddot{\theta} + 2\dot{r}\dot{\theta}) \hat{e}_\theta$$

Calculamos las derivadas necesarias:
\begin{itemize}
    \item $\dot{r} = -k\omega_0 r(t)$
    \item $\ddot{r} = -k\omega_0 \dot{r} = -k\omega_0 \cdot (-k\omega_0 r) = k^2 \omega_0^2 r(t)$
    \item $\dot{\theta} = \omega_0$
    \item $\ddot{\theta} = 0$
\end{itemize}

Componente radial:
$$a_r = \ddot{r} - r\dot{\theta}^2 = k^2 \omega_0^2 r - r \omega_0^2 = (k^2 - 1) \omega_0^2 r$$

Componente transversal:
$$a_\theta = r \cdot 0 + 2 \cdot (-k\omega_0 r) \cdot \omega_0 = -2k\omega_0^2 r$$

La aceleración total es:
\begin{equation}
\boxed{\vec{a} = (k^2 - 1) \omega_0^2 r(t) \, \hat{e}_r - 2k\omega_0^2 r(t) \, \hat{e}_\theta}
\end{equation}

o, factorizando:
\begin{equation}
\vec{a} = \omega_0^2 r(t) \left[ (k^2 - 1) \hat{e}_r - 2k \, \hat{e}_\theta \right]
\end{equation}

\paragraph{Análisis de la Aceleración}

\begin{itemize}
    \item Si $k < 1$: La componente radial apunta hacia afuera (positiva), aunque el punto se acerca al origen. Esto ocurre porque la velocidad radial disminuye en magnitud; el cambio de velocidad radial implica aceleración hacia afuera.
    \item Si $k = 1$: La componente radial es cero; solo existe aceleración transversal (similar a movimiento circular uniforme de radio variable).
    \item Si $k > 1$: La componente radial apunta hacia adentro, acelerando la convergencia.
    \item La componente transversal $a_\theta = -2k\omega_0^2 r$ siempre apunta en dirección negativa (contraria a la velocidad angular), ralentizando la rotación.
\end{itemize}

\subsubsection*{Pregunta 3: Rapidez Exponencial Decreciente}

La rapidez (magnitud de velocidad) es:
$$v = |\vec{v}| = \sqrt{(\dot{r})^2 + (r\dot{\theta})^2} = \sqrt{k^2\omega_0^2 r^2 + r^2 \omega_0^2}$$

$$v = \omega_0 r \sqrt{k^2 + 1} = \omega_0 \sqrt{k^2 + 1} \cdot r_0 e^{-k\omega_0 t}$$

\begin{equation}
\boxed{v(t) = v_0 e^{-k\omega_0 t}}
\end{equation}

donde $v_0 = \omega_0 r_0 \sqrt{k^2 + 1}$ es la rapidez inicial.

\paragraph{Demostración}

La rapidez disminuye exponencialmente con constante de tiempo $\tau = 1/(k\omega_0)$. Conforme $t \to \infty$:
- $r(t) \to 0$ (el punto converge al origen)
- $v(t) \to 0$ (la velocidad tiende a cero)
- El punto completa infinitas revoluciones alrededor del origen en tiempo finito, pero su velocidad disminuye exponencialmente.

Este comportamiento es típico de sistemas disipativos donde la energía cinética se reduce continuamente.
\end{solucion}
